\section{Discusión}
% Interpretacion propia de los datos

Como se mencionó en la sección~\ref{calculo_distancias}, el cálculo de distancias entre ciudades se realizó utilizando la fórmula de Haversine, que es adecuada para calcular 
distancias entre puntos en la superficie de una esfera, como la Tierra. Esta elección es crucial para obtener resultados precisos en el contexto del TSP, ya que las distancias 
reales entre ciudades afectan directamente la optimización de la ruta. Este es el primer punto a considerar al analizar los resultados obtenidos ya que el calculo con coordenadas
planas hubiera introducido errores significativos en la estimación de distancias reales entre ciudades ya que se consideraria una linea recta que en una esfera corta partes de la 
superficie, estas ligeras curvaturas pueden verse en los mapas de las rutas obtenidas.

Para la parte de la cruza, se utilizaron dos métodos: PMX y PBX, ambos diseñados para preservar la validez de las rutas generadas el principal componente en el PMX durante la
implementacion fue asegurar que no se repitieran ciudades en la ruta hija, lo cual se logró mediante un mapeo cuidadoso de los segmentos intercambiados entre los padres ya que en 
primeras instancias se presentaron errores donde se hacia un ciclo infinito por que no existian posiciones validas para las ciudades, por lo cual se implementaron validaciones 
adicionales para garantizar que cada ciudad apareciera exactamente una vez en la ruta resultante. En el caso del PBX, el desafío principal fue mantener el orden relativo de las 
ciudades no seleccionadas, lo cual se logró mediante la eliminación cuidadosa de duplicados y el llenado de las posiciones restantes con los elementos sobrantes del segundo padre.

En cuanto a los métodos de mutación, la mutación por barajado (scramble) y la mutación por intercambio de bloques (block transpose) fueron seleccionadas por su capacidad para
introducir diversidad en la población sin alterar significativamente la estructura de las rutas. La mutación por barajado altera el orden interno de un subconjunto de ciudades,
lo que puede ayudar a explorar nuevas combinaciones sin cambiar la posición relativa de las ciudades fuera del subconjunto afectado. Por otro lado, la mutación por intercambio de
bloques modifica la posición de segmentos enteros de la ruta, lo que puede resultar en cambios más drásticos en la estructura de la solución, pero también en una mayor exploración 
del espacio de búsqueda.

Con respecto a los resultados obtenidos, es evidente que el tamaño de la población inicial tiene un impacto significativo en la diversidad de soluciones exploradas por el algoritmo.
Una población más grande tiende a ofrecer una mayor variedad de rutas, lo que puede ayudar a evitar la convergencia prematura a óptimos locales. Sin embargo, también implica un 
mayor costo computacional, ya que se deben evaluar más soluciones en cada generación. Tambien podemos observar que la cruza PBX tiende a generar soluciones ligeramente menos óptimas 
pero más rapidas al momento de converger en comparación con PMX en este caso específico pero este ultimo presenta mejores resultados, aunque ambas cruces demostraron ser efectivas 
para explorar el espacio de soluciones.

Podemos observar que tambien el tener tantas generaciones (1000) puede ser contraproducente ya que en la mayoria de los casos la aptitud se estabiliza mucho antes de llegar a este 
punto, por lo que se podria considerar reducir el numero de generaciones para mejorar la eficiencia del algoritmo sin sacrificar significativamente o principalmente buscar un mejor
criterio de paro que no sea un numero fijo de generaciones.